\section{Scenarios}

\subsection*{THE SCENARIOS}

There are nineteen separate scenarios in the SHENANDOAH Game, each of which is
a game in itself, but each of which uses the same set of rules. Each scenario
is linked to the next by a running historical commentary, and each scenario
contains all the information reauired to set-up and play it. Many scenarios
contain special rules that apply only to that scenario (listed as "SPECIAL"). A
few conventions were used in writing out the scenarios, which are explained
below:

1. On the Initial Placement of the units, if only a hex coordinate is given,
the units must be placed directly in that hex. If the words "and adjacent"
appear next to the hex coordinate, the units may be spread in any desired
manner in the given hex and/or all hexes adjacent to it.

2. Garrison units are designated by being enclosed in [Brackets].

3. Advanced Game units, which would not be used in the Basic Game are
designated by being enclosed in (Parenthesis).

4, The listed Withdrawal Hexes are the only ones through which units can exit
the Mapboard to reduce their side's Resource Factor.

5. SHENANDOAH was play-tested during its development by a number of experienced
gamers. Their comments on each scenario are included to aide players trying a
scenario for the first time in formulating their strategies.

6. The Historical Result gives the results of the actual campaign based on the
point system used in this game.



\begin{scenariodata}
  SCENARIO SEQUENCE OF SETUP
  \begin{enumerate}
    \item First phase player sets up first.
    \item Second phase player sets up next.
    \item Partisans are always set up last. (Anywhere, over two hexes from nearest
  Union forces.)
  \end{enumerate}

  \begin{itemize}
    \item I = Infantry
    \item C = Cavalry
    \item A = Artillery
    \item H = Horse Artillery
    \item F = Fortification Counter
    \item S = Supply Counter
    \item W = Wagon Counter
  \end{itemize}
\end{scenariodata}

General Officer Counters are indicated by the general's name. Partisans are
indicated by the name of the Partisan (i.e. 10 | means ten combat factors of
infantry).

\subsection*{HINTS ON PLAY}

To play SHENANDOAH well, players need to organize and deploy their forces
effectively. Every player has his own "style," and will soon develop his own
favorite strategies on how best to play the game. The ideas presented here will
help new players in developing their own plans and strategies.

\textbf{GENERAL:} Wagons should be kept in separate stacks from the main body of the
army. Give them an adequate guard, and let these Stacks keep up as well as
possible; if they are kept with the large fighting stacks, they will hinder
their mobility. Form the infantry and artillery into a number of separate
Stacks, each large enough (ten or more Combat Factors) to be able to fight
effectively, take casualties, and keep on fighting. Form the cavalry and horse
artillery into smaller Stacks (three to six Combat Factors) to cover more
ground; avoid heavy combat with these units, but each Stack should be large
enough to engage in a Round or two of Combat without being annihilated, if
necessary.

Do not hesitate to spend a few points for additional wagon and/or supply
counters if they are needed. Keep in mind that there are two sides to the
Massanutten/Peaked Mountain range; regardless of which side the main effort is
on, the other side must always contain an ample security and delaying force.
Play to win each scenario. Merely playing to keep the other side from winning
results only in drawn games.

\textbf{CONFEDERATE:} Confederate wagons do not need to be as heavily escorted as Union
ones, as there are less Union opportunities to raid the Confederate rear areas.
Hidden movement abilities must be used to the utmost by feinting with "fake"
units in one direction while actually moving in another direction. Traps can be
set to draw Union attacks on strong forces that appear weak. The best method
for concealing movement is to set up a "cavalry screen" along the front -- this
screen consists of cavalry stacks only one or two Combat Factors each, and one
or two hexes apart (care must be taken to avoid having a part of the screen
surrounded). Although too weak to fight, such a screen has to be pushed back
and/or penetrated for the Union player to find the main Confederate army.

Superior leadership must be exploited; on the average, the Confederate ability
to force march and rally from disorders is far superior to the same
Union abilities. In most scenarios, the Rebels should avoid long costly
battles with large Union forces. Limit attacks to the weaker Union Stacks,
especially those with no General Officer. The Combat Tables greatly favor the
defender, and the Confederates should fight defensively most of the time;
attacking only when the odds greatly favor success.

The Partisan units are invaluable. They can operate near the B\&O to keep that
line cut by destroying unguarded stretches, and attacking weakly held
hexes (watch for unsupplied garrison units). Supported by some supplied
regular cavalry and horse artillery, the Partisans can threaten the
annihilation of any weak units along the B\&O line. Partisans are also useful
for operating in the rear areas of the Union field army, raiding wagon and
supply Stacks, cutting off the flow of fresh supplies, and in cutting off the
retreat of units attacked by the main Confederate forces. "Ambushes" can be
laid by being hidden adjacent to a rail or road hex down which an unescorted or
poorly-escorted supply unit may move during the Union Player's Phase.

\textbf{UNION:} Union wagon and supply Stacks must be well-guarded at all
times by infantry and/or artillery. Cavalry should be used for reconnaisance;
move them first, and move within two hexes of as many Confederate Stacks as
possible to separate the real ones from the "fake" Stacks. Once this is done,
the more powerful infantry forces will be needed to push back cavalry screens,
and to move adjacent to enemy units to determine their exact composition.

The biggest Union advantage in most scenarios is their overall superiority in
numbers, and the lesser Victory Point value of Union Combat Factors. The Union
Player should fight as often as possible (unless hopelessly outclassed), as
roughly even casualty rates ("attrition") give the Union more Victory Points
than the Confederates can gain.

In guarding the B\&O Railroad line, a decent sized force must be kept in
Harper's Ferry at all times. Spread the remainder of the B\&O garrison along
the rail line, guarding the bridges especially well. Fortify the most important
points, and bring in extra supply counters, if needed, to ensure that all
detachments are supplied. If Harper's Ferry is threatened by a large Rebel
force, get as much of the garrison over there as possible -- Harper's Ferry is
by far the most important hex along the B\&O. If Partisan units seem to be
hovering close to the B\&O line, conduct some "sweeps" within two hexes of the
track; such sweeps may be able to find, surround, and destroy them. In
Scenarios \#10 through \#16, the hexes adjacent to the track line could be
devastated; this earns no Victory Points, but greatly reduces the effectiveness
of the Partisans in raiding the line.

\subsection*{HISTORICAL SUMMARY}

On the evening of April 17, 1861, only three days after the surrender of Fort
Sumter, the tiny U.S. army garrison at Harper's Ferry fired the buildings of
the arsenal there, and retired across the Potomac. A force of Virginia militia
entered the town on the following morning. This was the curtain-riser on the
Civil War drama which would be acted out in the Shenandoah Valley and its
environs.

The Sheandoah Valley was an important, though secondary, theater of operations
for both sides during the war. Both sides desired to control the region, but
more pressing strategic considerations forced both sides to attempt to do so
with a minimum of forces.

Large forces were concentrated in the Valley in the early months of the war,
but there was little fighting. The Confederate Army of Sheandoah made a rapid
rail movement to join Beauregard at Manassas Junction to win the decisive First
Battle of Bull Run, and did not return to the Valley again.

By autumn of 1861 only small forces were again in the Sheandoah region. The
Confederates were commanded by Thomas J. Jackson (1824-1863), a graduate of
West Point, and a Mexican War veteran. One of the heroes of the Bull Run
battle, Jackson was regarded as either an eccentric or a madman, depending on
who you talked to. His first act on hearing of the secession of Virginia had
been to propose to his local pastor that all Christians should be called upon
to unite in prayer. His second act, being a true warrior from the pages of the
Old Testament, had been to offer his sword to his native state of Virginia. He
sucked on lemons, wore a disreputable uniform, rode a nag that most cavalry
privates would have been ashamed to be seen on, often preached to his men, and
seemed to fight all of his battles on Sundays. His sobriquet of "Stonewall" was
singularly inappropriate for a man who hated to fight defensively, and was
possibly the most aggressive commander to ever wear an American uniform. The
men he led more commonly referred to him as "Old Jack," or "Old Blue Light"
(due to the way his eyes would blaze during a battle).

The Rebel cavalry was led by Turner Ashby (1828-1862), an unconventional leader
who was idolized by his men (they would brag on the fact that he had
participated in twenty-eight different actions during the course of one month).
Other Confederate cavalry in the Valley consisted of numerous bands of
''Partisan Rangers," the most effective being led by "Handsome" Harry Gilmore,
and John ''Hanse" McNeill.

The Unicn troops in the Shenandoah were commanded by Nathaniel P. Banks
(1816-1894), an important Republican politician from Massa- chusetts with few
military qualifications. His subordinate, James Shields (1810-1879) was an
Irish-born, Oxford educated immigrant. Shields had served three times in the
United States Senate, each time representing a different state (Illinois,
Minnesota, and California). He had been a Brigadier Genera! of Illinois
volunteers during the Mexican War, and had once almost fought a duel with
Abraham Lincoln. A combative leader, Sheilds was more aggressive than he was
able.

Banks had a much larger force than Jackson, and he had orders to drive the
Confederates from Winchester to stop the constant attacks on the B\&O Railroad.
Following this advance, a sufficient force was to be left to-contain further
Confederate excursions, and Banks was to rejoin the main army for the push on
Richmond. By the middle of March, 1862, Winchester was in Federal hands, and
Banks was marching out of the Valley with one of his divisions, leaving Shields
to guard their gains. At this point, Jackson, having orders to prevent just
such a withdrawal of Union troops from the Shenandoah, decided to strike...
